\chapter{Glossary}
\label{app:glossary}

\begin{description}[leftmargin=1.6em,font=\bfseries,itemsep=2pt]
  \item[Action] The policy's output for one step; here two numbers in $[-1,1]$.
  \item[Actuator (implicit)] A PD drive computed inside the physics solver from
    position/velocity targets and gains. \code{ImplicitActuatorCfg}.
  \item[Advantage] How much better an action turned out than the critic
    expected. Estimated with GAE.
  \item[Articulation] A tree of rigid bodies connected by joints, with one
    root. The car. \code{Articulation} / \code{ArticulationCfg}.
  \item[Asset] Anything the scene spawns and wraps in a Python object:
    articulation, rigid object, MPM object, or a plain \code{AssetBaseCfg}.
  \item[Checkpoint] A saved snapshot of the networks, \code{model_N.pt}.
  \item[configclass] Isaac Lab's dataclass wrapper used for every config.
  \item[Coupler] Newton's experimental solver that runs several solvers on
    disjoint parts of one model and links them. \code{CouplerProxyCfg}.
  \item[Coupler entry] One named sub-solver and the bodies/particles it owns.
    \code{CouplerEntryCfg}.
  \item[Critic / value function] The network that predicts expected return
    from an observation.
  \item[CUDA graph] A recorded sequence of GPU kernels replayed without CPU
    overhead; requires fixed buffer sizes.
  \item[Decimation] Physics steps per control step. 2 here.
  \item[DirectRLEnv] Isaac Lab's base class for environments written as
    explicit methods (as opposed to the manager-based workflow).
  \item[Entry point] Installed package metadata that Isaac Lab uses to
    discover task packages (\code{isaaclab.tasks} group).
  \item[Env / environment instance] One copy of the scene on the env grid.
  \item[Env origin] World position of an env's local frame.
    \code{scene.env_origins}.
  \item[Episode] A run from reset to done; 2~s or an early termination.
  \item[GAE] Generalised Advantage Estimation, parameter $\lambda$.
  \item[Hydra] The config-override system behind \code{physics=NAME} and
    \code{env.xxx=value} command-line arguments.
  \item[InteractiveScene] Turns a scene cfg into spawned prims and asset
    objects; owns \code{env_origins}.
  \item[Iteration] One PPO cycle: rollout, advantage, update.
  \item[Kinematic body] A rigid body with a pose but no dynamics; the slab.
  \item[Lagged coupling] Proxy poses sent forward each step, feedback
    impulses applied one step late.
  \item[MJCF] MuJoCo's XML model format. \code{tricycle.py} holds one.
  \item[MJWarp] MuJoCo re-implemented in Warp; Newton's rigid-body solver.
  \item[MPM] Material Point Method: particles plus a background grid for
    continua such as soil.
  \item[MPMObject] Isaac Lab asset wrapping one particle cloud per env.
  \item[Newton] NVIDIA's GPU physics engine on Warp, Isaac Lab's default
    backend on develop.
  \item[Observation] The policy's input; 15 numbers here.
  \item[Policy / actor] The network mapping observations to an action
    distribution.
  \item[PPO] Proximal Policy Optimisation, the learning algorithm.
  \item[Prim] A node in a USD stage. Prim paths look like
    \srcpath{/World/envs/env_0/Tricycle}.
  \item[Proxy] A ghost copy of a body from one coupler entry placed inside
    another so they can interact.
  \item[Rigid object] A single rigid body asset. \code{RigidObjectCfg}.
  \item[Rollout] The batch of experience collected before an update.
  \item[Root state] Position, quaternion, linear and angular velocity of an
    articulation's root body: 13 numbers.
  \item[rsl\_rl] The learning library (ETH Zurich Robotic Systems Lab).
  \item[SimulationContext] The Isaac Lab object that owns the stage and the
    physics manager.
  \item[Sparse grid] MPM grid that allocates cells only near particles; sized
    by the four capacity caps.
  \item[Time-out] Episode cut by the clock, not by failure; bootstrapped by
    the learner.
  \item[USD] Universal Scene Description, the scene file format Isaac Lab
    builds physics from.
  \item[Warp] NVIDIA's Python-to-CUDA kernel language that Newton and Isaac
    Lab are written in.
\end{description}

\chapter{What changed in this revision}
\label{app:changes}

All changes are on the branch named on the title page. None changes the
task's observation, action, reward or termination definitions.

\section*{The car body}
\begin{itemize}[leftmargin=1.6em]
  \item \code{tricycle.py}: the single box chassis became a triangular frame
    (two rails, crossbar, deck). Wheel positions, radii, masses, joint names
    and body names are unchanged; the chassis total mass is still 5.0~kg.
    Geometry is computed from a few named constants and a \code{_rail}
    helper. The import-time \code{write_mjcf()} call was removed.
  \item \srcpath{assets/tricycle/}: the committed USD was regenerated to match.
    The four new \code{Cube} prims replace \code{chassis_geom} in
    \code{base.usda}, \code{physics.usda} and \code{mujoco.usda}. The edit
    was validated by composing the stage with OpenUSD: every prim carries the
    same schemas as before and the four masses sum to 5.0~kg. The new MJCF
    was also compiled with MuJoCo and driven for several seconds on a plane
    (stands, drives straight, turns without tipping).
  \item If anything ever looks wrong, the converter is the source of truth:
    rerun the two conversion commands.
\end{itemize}

\section*{Code clean-up (behaviour preserved)}
\begin{itemize}[leftmargin=1.6em]
  \item \code{tricycle_env_cfg.py}: the hard-coded Windows asset path became
    a path relative to the file; stale comments about \code{MjcfFileCfg}
    were removed; the strip comment now states the real particle count.
  \item \code{tricycle_env.py}: removed a stale \code{VERIFY} marker on an
    import; type annotations on \code{self.car}/\code{self.soil}; the
    unreachable \code{_ALL_INDICES} fallback replaced by \code{torch.arange};
    the \code{soil.reset} comment now states what the call does.
  \item \srcpath{tricycle/__init__.py}: header shows the current CLI commands
    instead of the removed \code{train.py} script and \code{-v0} id.
  \item \srcpath{agents/rsl_rl_ppo_cfg.py}: the deprecated runner-level
    \srcpath{empirical_normalization} was replaced by the two policy-level
    normalisation flags it used to set.
  \item Deleted \srcpath{tricycle/soil.py}, an unused copy of
    \srcpath{excavation/soil.py}; deleted the stray root file
    \code{__init__ .py}.
  \item Added \code{.gitignore}; \srcpath{__pycache__/} and
    \srcpath{luna_hifi_tasks.egg-info/} are no longer tracked (they are
    regenerated by Python and pip).
\end{itemize}

\section*{Documentation}
\begin{itemize}[leftmargin=1.6em]
  \item \code{RUNBOOK.md}: a note on the body change under section 5.
  \item \srcpath{docs/tricycle_manual/}: this document and its LaTeX source.
    Rebuild with \srcpath{./build.sh} (needs TeX Live with \code{latexmk},
    \code{tcolorbox}, \code{tikz}, \code{listings}).
\end{itemize}

\section*{Pulling these changes}
If \code{git checkout} or \code{git pull} complains that local changes to
\code{.pyc} or \code{egg-info} files would be overwritten, those are the
formerly tracked generated files; discard them with \code{git checkout -- .}
(or \code{git stash}) and retry.

No reinstall is needed: \code{pyproject.toml} is unchanged by this revision,
and the \srcpath{egg-info} the pull removes is a build leftover rather than
the metadata the CLI reads (Chapter~\ref{ch:pyproject}). \code{tricycle.py}
did change, so reconvert the asset:
\begin{lstlisting}[style=shell]
.\isaaclab.bat -p -c "from luna_hifi_tasks.tricycle.tricycle import write_mjcf; print(write_mjcf())"
.\isaaclab.bat -p scripts\tools\convert_mjcf.py C:\Users\hanse\AppData\Local\Temp\luna_hifi_assets\tricycle.xml C:\Users\hanse\Documents\Luna_HiFi\assets\tricycle.usd
\end{lstlisting}
If the registry check ever comes back empty, reinstall once and it will be
populated again:
\begin{lstlisting}[style=shell]
.\isaaclab.bat -p -m pip install -e C:\Users\hanse\Documents\Luna_HiFi
\end{lstlisting}

\chapter{Tensor shapes cheat sheet}
\label{app:shapes}

$N$ is the number of envs.

\begin{center}
\small
\begin{tabular}{@{}lll@{}}
\toprule
Expression & Shape & Meaning \\
\midrule
\code{actions} (to \code{_pre_physics_step}) & $(N,2)$ & throttle, steer in $[-1,1]$ \\
\code{self._actions} & $(N,2)$ & clamped copy, also observed \\
\code{car.data.root_pos_w} & $(N,3)$ & root position, world frame \\
\srcpath{car.data.root_lin_vel_b} & $(N,3)$ & root linear velocity, body frame \\
\srcpath{car.data.root_ang_vel_b} & $(N,3)$ & root angular velocity, body frame \\
\srcpath{car.data.projected_gravity_b} & $(N,3)$ & gravity direction in body frame \\
\code{car.data.joint_pos} & $(N,J)$ & all joint positions; $J=4$ here \\
\code{car.data.joint_vel} & $(N,J)$ & all joint velocities \\
\srcpath{car.data.default_root_state} & $(N,13)$ & pos(3) quat(4) lin vel(3) ang vel(3) \\
\srcpath{car.data.default_joint_pos} & $(N,J)$ & from the cfg \code{init_state} \\
\code{scene.env_origins} & $(N,3)$ & env origins, world frame \\
\code{obs["policy"]} & $(N,15)$ & the observation \\
reward & $(N,)$ & float \\
terminated, time\_out & $(N,)$ & bool \\
\code{episode_length_buf} & $(N,)$ & long, steps since reset \\
\srcpath{soil.data.nodal_pos_w} & $(N, 720, 3)$ & particle positions per env \\
\bottomrule
\end{tabular}
\end{center}

The joint count $J$ is 4: \code{steer}, \code{front_wheel},
\code{rear_left}, \code{rear_right}; the chassis free joint is the root
state, not a joint in this table. Their order in the tensors comes from the
USD and is what \code{find_joints} resolves.

\chapter{Where to look in the Isaac Lab source}
\label{app:sources}

Paths are relative to the Isaac Lab checkout (\code{develop}).

\begin{center}
\small
\begin{tabular}{@{}p{5.6cm}p{9.2cm}@{}}
\toprule
Concept & File \\
\midrule
CLI, entry-point loading & \srcpath{source/isaaclab/isaaclab/cli/__init__.py} \\
train / play backends & \srcpath{source/isaaclab_rl/isaaclab_rl/entrypoints/backends/train_rsl_rl.py}, \code{play_rsl_rl.py} \\
common CLI args, \code{--num_envs} override & \srcpath{source/isaaclab_rl/isaaclab_rl/entrypoints/common.py} \\
Hydra presets, \code{play_mode()} call & \srcpath{source/isaaclab_tasks/isaaclab_tasks/utils/hydra.py} \\
\code{DirectRLEnv} step/reset & \srcpath{source/isaaclab/isaaclab/envs/direct_rl_env.py} \\
\code{DirectRLEnvCfg} fields, base \code{play_mode} & \srcpath{source/isaaclab/isaaclab/envs/direct_rl_env_cfg.py} \\
\code{configclass} & \srcpath{source/isaaclab/isaaclab/utils/configclass.py} \\
spaces from ints & \srcpath{source/isaaclab/isaaclab/envs/utils/spaces.py} \\
\code{InteractiveScene} & \srcpath{source/isaaclab/isaaclab/scene/interactive_scene.py} \\
\code{SimulationCfg} & \srcpath{source/isaaclab/isaaclab/sim/simulation_cfg.py} \\
articulation API & \srcpath{source/isaaclab/isaaclab/assets/articulation/base_articulation.py} \\
root/joint data fields & \srcpath{source/isaaclab/isaaclab/assets/rigid_object/base_rigid_object_data.py} \\
actuator cfgs & \srcpath{source/isaaclab/isaaclab/actuators/actuator_base_cfg.py}, \code{actuator_pd_cfg.py} \\
\code{NewtonCfg} & \srcpath{source/isaaclab_newton/isaaclab_newton/physics/newton_manager_cfg.py} \\
\code{MJWarpSolverCfg} & \srcpath{source/isaaclab_newton/isaaclab_newton/physics/mjwarp_manager_cfg.py} \\
\code{MPMSolverCfg} & \srcpath{source/isaaclab_newton/isaaclab_newton/physics/mpm_manager_cfg.py} \\
MPM manager, \code{reset_solver_state} & \srcpath{source/isaaclab_newton/isaaclab_newton/physics/mpm_manager.py} \\
collision pipeline cfg & \srcpath{source/isaaclab_newton/isaaclab_newton/physics/newton_collision_cfg.py} \\
\code{MPMObject}, \code{reset} & \srcpath{source/isaaclab_newton/isaaclab_newton/assets/mpm_object/mpm_object.py} \\
\code{MPMGridCfg}, material & \srcpath{source/isaaclab_newton/isaaclab_newton/sim/spawners/mpm/mpm_cfg.py}, \code{mpm.py} \\
Newton collision props & \srcpath{source/isaaclab_newton/isaaclab_newton/sim/schemas/schemas_cfg.py} \\
coupler cfgs and manager & \srcpath{source/isaaclab_contrib/isaaclab_contrib/coupling/coupler_cfg.py}, \code{coupler.py} \\
reference coupled task & \srcpath{source/isaaclab_tasks/isaaclab_tasks/contrib/ur10_particle_push/ur10_particle_push_env_cfg.py} \\
Newton GL visualiser cfg & \srcpath{source/isaaclab_visualizers/isaaclab_visualizers/newton/newton_visualizer_cfg.py} \\
rsl\_rl cfg classes & \srcpath{source/isaaclab_rl/isaaclab_rl/rsl_rl/rl_cfg.py} \\
rsl\_rl cfg conversion & \srcpath{source/isaaclab_rl/isaaclab_rl/rsl_rl/utils.py} \\
rsl\_rl wrapper & \srcpath{source/isaaclab_rl/isaaclab_rl/rsl_rl/vecenv_wrapper.py} \\
MJCF converter & \srcpath{source/isaaclab/isaaclab/sim/converters/mjcf_converter.py}, \srcpath{scripts/tools/convert_mjcf.py} \\
\bottomrule
\end{tabular}
\end{center}

In the Newton package (\srcpath{newton/_src/}): the proxy coupler is
\srcpath{solvers/coupled/solver_coupled_proxy.py}, the implicit MPM solver is
\srcpath{solvers/implicit_mpm/solver_implicit_mpm.py}. In \code{rsl_rl}:
\srcpath{runners/on_policy_runner.py} and \srcpath{algorithms/ppo.py}.

\chapter{Numbers at a glance}
\label{app:numbers}

\section*{The car (chassis frame, metres, kilograms)}
\begin{center}
\small
\begin{tabular}{@{}llll@{}}
\toprule
Part & Centre $(x,y,z)$ & Half-sizes / radius & Mass \\
\midrule
left rail & $(0.035, 0.065, 0)$, yaw $-19.4^\circ$ & $(0.1961, 0.02, 0.015)$ & 0.8 \\
right rail & $(0.035, -0.065, 0)$, yaw $+19.4^\circ$ & $(0.1961, 0.02, 0.015)$ & 0.8 \\
crossbar & $(-0.15, 0, 0)$ & $(0.02, 0.15, 0.015)$ & 0.8 \\
deck & $(-0.08, 0, 0)$ & $(0.06, 0.065, 0.01)$ & 2.6 \\
fork capsule & $(0.22, 0, -0.04)$ & $r=0.012$, half-length 0.02 & 0.1 \\
front wheel & $(0.22, 0, -0.04)$ & $r=0.10$ & 0.5 \\
rear left wheel & $(-0.15, 0.16, -0.04)$ & $r=0.10$ & 0.5 \\
rear right wheel & $(-0.15, -0.16, -0.04)$ & $r=0.10$ & 0.5 \\
\midrule
chassis body & COM $(-0.054, 0, 0)$ & inertia $(0.062, 0.044, 0.019)$ & 5.0 \\
whole car & & wheelbase 0.37, track 0.32 & 6.6 \\
\bottomrule
\end{tabular}
\end{center}

\section*{The soil}
\begin{center}
\small
\begin{tabular}{@{}ll@{}}
\toprule
strip box (env-local) & $x\in[-0.4, 1.1]$, $y\in[-0.3, 0.3]$, $z\in[0.025, 0.085]$ \\
voxel & 0.05 m \\
particles per env & $30\times12\times2 = 720$, one per cell centre \\
particle mass & $0.05^3\times1800 = 0.225$ kg \\
material & $\rho=1800$, friction $0.84$, yield pressure $10^{12}$ Pa, cohesion 0 \\
slab & $2.5\times1.6\times0.1$ m centred at $(0.35, 0, -0.05)$, kinematic, invisible \\
grid caps (active/leaf/lower/upper) & $16384 / 8192 / 4096 / 1024$, totals for all envs \\
\bottomrule
\end{tabular}
\end{center}

\section*{Time}
\begin{center}
\small
\begin{tabular}{@{}ll@{}}
\toprule
physics step & 10 ms (rigid entry: 2 sub-steps of 5 ms; MPM: 1) \\
control step & 20 ms (50 Hz) \\
episode & 2 s $=$ 100 control steps $=$ 200 physics steps \\
rollout & 50 control steps per env per iteration \\
training & 300 iterations, checkpoint every 50 \\
\bottomrule
\end{tabular}
\end{center}
